% Capitulo: De usuários a tokens: custo e consumo
\chapter{De usuários a tokens: custo e consumo}
\label{cap:custo-consumo}

\rascunho{
  \textbf{Ideia do capítulo.} As ferramentas tradicionais de dados eram cobradas por
  usuário: o custo era fixo, e o gerente controlava quem tinha acesso. A tendência agora
  é cobrar por consumo (créditos de processamento, tokens de modelos de linguagem), e a
  gestão desse consumo virou preocupação de todo gerente de dados, para controlar custo
  e recurso.
  \begin{itemize}
    \item A mudança começou com os armazéns de dados na nuvem, que cobram pelo
      processamento (ver o capítulo~\ref{cap:nuvem}), e se acelerou com a IA, em que
      cada pergunta tem um custo em tokens.
    \item Ferramentas de BI também estão migrando para modelos de crédito ou capacidade.
      \pendente{conferir exemplos atuais de ferramentas com modelo por consumo, inclusive as que eu uso}
    \item Liga direto com o capítulo~\ref{cap:frequencia}: a frequência de atualização
      é uma das maiores alavancas de consumo.
    \item Referência para a disciplina de gestão de custo em nuvem (FinOps):
      \textcite{storment2023}.
  \end{itemize}
}

% ---------------------------------------------------------------------------
\section{O modelo por usuário}
\label{sec:modelo-usuario}

\rascunho{
  \begin{itemize}
    \item Licença por pessoa, muitas vezes em níveis (quem cria, quem explora, quem só
      visualiza). Custo previsível, negociado na renovação do contrato.
    \item A alavanca do gerente era o acesso: quem precisa de licença e de qual nível.
    \item O desperdício típico: licença paga para quem não entra há meses.
  \end{itemize}
}

\pendente{escrever}

% ---------------------------------------------------------------------------
\section{Créditos, consumo e tokens}
\label{sec:modelo-consumo}

\rascunho{
  \begin{itemize}
    \item No modelo por consumo, paga-se pelo que foi processado: consultas, atualizações,
      horas de computação, tokens de entrada e de saída de um modelo de linguagem.
    \item A Tabela~\ref{tab:usuario-consumo} compara os dois modelos.
    \item Vantagem: paga-se pelo uso real, e um time pequeno começa barato. Risco: o
      custo fica variável, e um processo mal configurado pode consumir o orçamento do
      mês em poucos dias.
  \end{itemize}
}

\pendente{escrever}

\begin{table}[htbp]
  \centering
  \caption{Cobrança por usuário e por consumo.}
  \label{tab:usuario-consumo}
  \small
  \begin{tabularx}{\textwidth}{@{}>{\raggedright\arraybackslash}p{0.2\textwidth}
      >{\raggedright\arraybackslash}X>{\raggedright\arraybackslash}X@{}}
    \toprule
    & Por usuário & Por consumo (créditos, tokens) \\
    \midrule
    O que se paga & Uma licença por pessoa com acesso & O que foi processado: consultas, atualizações, tokens \\
    \addlinespace
    Previsibilidade & Alta, custo fixo por período & Baixa, varia com o uso \\
    \addlinespace
    Alavanca de controle & Quem tem acesso e com qual nível & O que roda, com que frequência e com qual modelo \\
    \addlinespace
    Desperdício típico & Licença de quem não usa & Atualização ou consulta que ninguém aproveita \\
    \addlinespace
    Quem acompanha & Compras e TI, na renovação & O gerente do time, de forma contínua \\
    \bottomrule
  \end{tabularx}
  \fonte{elaborado pelo autor.}
\end{table}

\antesdepois{O custo da ferramenta era uma linha fixa no orçamento. O gerente controlava
  quem tinha licença, e o uso, alto ou baixo, não mudava a fatura.}
{O custo acompanha o uso. O gerente controla o que roda, com que frequência e com qual
  modelo, e acompanha o consumo como acompanha qualquer outro indicador do time.}

% ---------------------------------------------------------------------------
\section{O gerente como gestor de consumo}
\label{sec:gestao-consumo}

\rascunho{
  \begin{itemize}
    \item Orçamento de créditos ou tokens por time e por produto, com alertas de consumo
      antes de estourar.
    \item Custo por produto de dados (dashboard, pipeline, agente), usando as métricas de
      metadados do capítulo~\ref{cap:frequencia} (\veja{sec:metricas-uso}).
    \item Mostrar o custo para quem pede: quando o stakeholder vê quanto custa atualizar
      de hora em hora, a conversa sobre frequência muda.
    \item As práticas de FinOps: visibilidade, responsabilidade de cada time pelo próprio
      custo, otimização contínua \parencite{storment2023}.
    \item Relato: como eu acompanho consumo hoje. \pendente{escrever com relato, sem nomes nem valores exatos}
  \end{itemize}
}

\pendente{escrever}

% ---------------------------------------------------------------------------
\section{Tokens no trabalho com dados}
\label{sec:tokens-dados}

\rascunho{
  \begin{itemize}
    \item Onde os tokens aparecem: assistentes de SQL, perguntas em linguagem natural
      sobre dashboards, agentes que rodam tarefas inteiras, documentação gerada.
    \item O que faz o custo crescer: contexto grande (mandar a tabela inteira para o
      modelo), perguntas repetidas sem cache, usar o modelo mais caro para tudo.
    \item O que ajuda: pré-calcular o que é perguntado sempre, escolher o modelo pelo
      tamanho da tarefa, limitar o contexto ao necessário, medir tokens por resposta
      útil.
    \item A mesma lógica do capítulo anterior: não adianta responder em segundos uma
      pergunta que poderia vir num relatório semanal.
  \end{itemize}
}

\pendente{escrever}

% ---------------------------------------------------------------------------
\begin{leituras}
  \leitura{storment2023}{A referência sobre FinOps: como dar visibilidade ao custo de
    nuvem, distribuir a responsabilidade entre os times e otimizar sem travar o trabalho.}
\end{leituras}
